\section{From normal form sentences to general sentences} \label{app:normal}

In the main body all our results are proven for sentences in normal form.
In this appendix we will show that they also hold for general sentences.
In particular we will sketch the algorithms that transform an arbitrary ($\GPtwo$/$\Ctwo$) sentence into another sentence in the normal form and recover the properties of the original input sentence from the constructed normal form sentence.
We emphasize that the transformations we provide are a routine adaptation of the textbook transformation to Scott normal form: one can find similar results 
in~\cite{twovarpres}, for the $\GPtwo$ case,
and in~\cite{ctwocomplexity} for the $\Ctwo$ case.

We start with $\GPtwo$ sentences, showing that we can find an equisatisfiable formula in normal form.
Note that in the normal form we allow double universal quantifiers of the form:
$$
\forall x \forall y\ \left(R(x,y) \land x \neq y\right) \to \alpha(x, y)
$$
which can be considered as a shorthand for the  sentence:
$$
\forall x 
\left(
\presby{R(x, y) \land x \neq y\land \neg \alpha(x,y)}
\ =\ 0 
\right)
$$

Recall that the normal form of a $\GPtwo$ sentence has the shape:
\begin{equation*}
    \varphi\ :=\ 
    \forall x\ \gamma(x)\ \land\ 
    \bigwedge_{i \in \intsinterval{m}} \forall x \forall y\ \left(R_i(x,y) \land x \neq y\right) \to \alpha_i(x, y)\ \land\ 
    \bigwedge_{i \in \intsinterval{n}} \forall x\ U_i(x) \to P_i(x),
\end{equation*}
where
\begin{compactitem}
    \item $\gamma(x)$ is a quantifier-free formula,
    \item each $\alpha_i(x, y)$ is a quantifier-free formula,
    \item each $P_i(x)$ is a Presburger quantified formula of the form: 
    \begin{equation*}
        P_i(x)\ :=\ \left(
        \sum_{t \in \intsinterval{m}} \lambda_{i, t} \cdot \presby{R_t(x, y) \land x \neq y}
        \ \circledast_i\ \delta_i \right).
    \end{equation*}
\end{compactitem}

The main idea is as follows.
Let $\psi$ be an arbitrary $\GPtwo$ sentence.
By pushing every negation inward, we may assume that it is in negation normal form,
i.e., every negation is applied only on atomic formulas.
We then repeatedly replace every subformula $\nu(x)$ of one free variable $x$ with a fresh unary predicate $U(x)$, and conjoin the original formula with the ``definition'' of $U(x)$:
$$
\forall x \ U(x) \to \nu(x)
$$
Intuitively what this means is that
the subformula $\nu(x)$ is renamed as a unary predicate $U(x)$.
The algorithm starts bottom-up starting from the subformula $\nu(x)$
with the lowest quantifier rank.
It can be shown that the number of renaming required to obtain the normal form is linear in the length of the original sentence $\psi$~\cite[Lemma 3.3]{twovarpres} and the whole transformation runs in linear time.

Note that unguarded unary quantification is allowed in $\GPtwo$.
For universally-quantified formulas, $\forall x\ U(x)$ is already captured by the normal form.
For existential-quantified $\psi := \exists x\ U(x)$,
we introduce fresh unary predicates $U'$ and binary predicate $R_U$, and replace $\psi$ by 
\begin{equation*}
    \left(
        \forall x\ U(x) \lor U'(x)
    \right)
    \ \land\ 
    \left(
        \forall x\ U'(x) \to \presby{R_U(x, y) \land U(y)} = 1
    \right).
\end{equation*}

Note also that we can recover the model of the original sentence from the model of the constructed normal form sentence.
We simply project out the fresh unary predicates introduced when constructing the normal form sentence. 

Now, we turn to $\Ctwo$ sentences, and recall from the body of the paper the normal form we utilized:
\begin{equation*}
    \varphi\ :=\ 
    \forall x\ \gamma(x)\ \land\ 
    \forall x\ \forall y\ \left(x \neq y \to \alpha(x, y)\right)\ \land\ 
    \bigwedge_{i \in \intsinterval{m'}} \forall x\ \exists^{=k_i} y\ \left(R_i(x,y) \land x \neq y\right),
\end{equation*}
where $\gamma(x)$ and $\alpha(x, y)$ are quantifier-free formulas,
$R_i$ are binary predicates,
and $k_i \in \bbN$.
 
The main idea of the transformation is similar to the $\GPtwo$ case:
Rename every subformula $\nu(x)$ of one free variable $x$ with a fresh unary predicate $U(x)$ and conjunct the original sentence with $\forall x\  U(x)\to \nu(x)$.
It was shown in~\cite{ctwocomplexity,ianrevisited} that there is a linear time algorithm that transforms
an arbitrary $\Ctwo$ sentence $\varphi$ over vocabulary $\vocab$ to a sentence $\varphi^*$ in normal form (with additional unary predicates $V_1,\ldots,V_p$) such that:
\begin{compactitem}
\item 
For every model $\cA\models \varphi$ with domain size at least $1+\max_{i\in[m']}k_i$,
there is a model $\cA^*\models \varphi^*$, where $\cA^*$ is an extension of $\cA$ to the vocabulary $\vocab^*:=\vocab\cup\{V_1,\ldots,V_p\}$ with
$A=A^*$ and $S^{\cA}=S^{\cA^*}$, for every relation symbol $S\in\vocab$.
\item
Conversely,
for every model $\cA^*\models \varphi^*$ with domain size at least $1+\max_{i\in[m']}k_i$,
there is a model $\cA\models \varphi$, where $\cA$ is obtained by projecting out the predicates $V_1,\ldots,V_p$ in $\cA^*$.
\end{compactitem}
Using these two properties, we can generalise Theorem~\ref{thm:spectrasemilinear} to arbitrary $\Ctwo$ sentences.
First, all the sizes of the models up to 
$1+\max_{i\in [m']} k_i$ can be encoded inside the corresponding Presburger formula $\Psi^\varphi(\bfx)$
mentioned in Lemma~\ref{thm:c2_spec}.
Second, we can recover $\pispec$ from $\Pi^*\textsf{-SPEC}$,
where $\Pi$ and $\Pi^*$ are the set of all 1-types w.r.t. $\vocab$ and $\vocab^*$,
via the following identity, for every 1-type $\pi$ over $\vocab$:
$$
x_{\pi} \ = \
\sum_{\pi^*\supseteq \pi \ \text{and}\ \pi^*\ \text{is 1-type over}\ \vocab^*} x_{\pi^*}
$$
